\documentclass[11pt]{article}
\usepackage[a4paper,margin=25mm]{geometry}
\usepackage{amsmath,amssymb,amsthm,mathtools,hyperref}
\newtheorem{theorem}{Theorem}\newtheorem{lemma}[theorem]{Lemma}
\newcommand{\R}{\mathbb R}\newcommand{\C}{\mathbb C}
\newcommand{\PP}{\mathcal P}\newcommand{\BB}{\mathcal B}
\newcommand{\dd}{\,d}\newcommand{\Gram}{\operatorname{Gram}}
\title{The absolute Gamma baseline and the complete action at logarithmic precision}
\author{Split-Zero research programme: mathematical continuation}
\date{19 September 2026}
\begin{document}\maketitle

\noindent The incoming note \texttt{07\_ABSOLUTE\_GAMMA\_BASELINE.md}
identified the previously unevaluated absolute Gamma coefficient. This proof
retains its complete root polynomial, original monic frames and source mass.
It strengthens its remainder to $O_h(\log q)$ for every fixed multiplicity.
The original four blocks are expressed using a single fixed potential and
three analytic factors, and a finite convexity argument controls the remaining
bounded root parameter. The previously proved arithmetic correction is then
returned to the absolute baseline and the complete canonical action.
The human analytic input is Borot--Guionnet \cite{BG}; it is applied after
a comparison of the complete quadratic forms on both complementary regions.

\section{Original objects and the determinant map}
Fix $0<\delta<1/2$, $\gamma>2$ and $m_0\in\mathbb N$. For the arithmetic
receivers these are the stipulated complete quartet data of the preceding
programme, with $h(s)$ its actual quartet factor in $2\xi(s)$.
The Gamma calculation itself requires only these fixed parameters. Put
\begin{equation}\tag{AG1}
 k=4\ell+1,\quad e=1+k(m_0-1),\quad q=e(k+1)^2,\quad c=k/2,
 \qquad
 \chi(S)=\prod_{a,b=0}^k
 [S-c-(2a-k)\delta-i(2b-k)\gamma]^e.
\end{equation}
The monic polynomial $Q(y)=i^{-q}\chi(c+iy)$ has its full roots
\begin{equation}\tag{AG2}
 \omega_{ab}=(2b-k)\gamma-i(2a-k)\delta
 \quad\hbox{with multiplicity }e.
\end{equation}
Conjugation and negation preserve the root multiset. Since $k$ is odd,
no root is real. Thus $Q$ is real, even and strictly positive on $\R$.
Here $q$ is divisible by four. Retain
\[
 \sigma(y)=\frac{|\Gamma(1/4+iy/2)|^2}{2\pi},\qquad
 \int_\R\sigma(y)\dd y=\sqrt{2\pi},\qquad \alpha=\pi/2.
\]
Let $H_N$ be the monic source Gram through degree $N$, $J_N$ the remainder
map into $\C[y]/Q$, and $G_N=(J_NH_N^{-1}J_N^*)^{-1}$ for $N\ge q-1$.
Its definition is the squared minimum norm in each original remainder fibre.
The monic source norms are
\begin{equation}\tag{AG3}
 \gamma_n=\sqrt{2\pi}\,n!(1/2)_n
 =\sqrt{2\pi}\frac{\Gamma(2n+1)}{4^n}.
\end{equation}
These are the original Meixner--Pollaczek specialization; the generating
function is $(1+t^2)^{-1/4}\exp(y\arctan t)$, with coefficients $p_n(y)/n!$.
The general polynomial convention is recorded in DLMF \cite{DLMF}.
Put $R_n=\det[\int y^{i+j}|Q|^2\sigma\dd y]_{0\le i,j<n}$ and $R_0=1$.
The frame $(1,\ldots,y^{q-1},Q,yQ,\ldots,y^{N-q}Q)$ has determinant one.
Taking its Schur complement, with every source/relation cross term present,
proves $\det G_N=\det H_N/R_{N-q+1}$. Consequently
\begin{equation}\tag{AG4}
 \BB_k[\sigma]=\log R_q+\log R_{q+1}-\log R_1
       -\sum_{n=q}^{2q}c_{n-q}\log\gamma_n,
 \quad c_0=c_q=1,\quad c_r=2\ (0<r<q).
\end{equation}
The affine substitution $S=c+iy$ has triangular coefficient matrix with
diagonal $1,i,\ldots,i^N$, of determinant modulus one. Its inverse is
$y=(S-c)/i$; it intertwines the relation injections and remainder maps.
Thus (AG4) computes the original $S$-coordinate return, without a new mass.

\section{The fixed equilibrium measure and its inverse moment}
Use the EIQ measure for $V(x)=\pi\sqrt x-2\log x$, with support
$[u^2,v^2]$. In the original endpoint convention
\begin{equation}\tag{AG5}
 \kappa^2=1-u^2/v^2,\qquad uK(\kappa)=2,\qquad vE(\kappa)=4,
 \quad F=-\min\mathcal E_{2,\pi},\quad
 L=\int\log x\dd\rho(x),\quad M=\int x^{-1}\dd\rho(x).
\end{equation}
All integrals use the original probability density, rather than a scaled
endpoint model. Its complete proof is EIQ1--26 in
\texttt{EIQ\_GEL\_COMPLETE.tex}. For application of the analytic theorem
we recall its decisive exact formulas. For general $a=u^2,b=v^2$,
$R_s=\sqrt{(a+s)(b+s)}$, and parameters $A,B>0$, they are
\begin{equation}\tag{AG6}
 h_{A,B}(x)=\frac{B}{2\pi x}\int_0^\infty
 \frac{\sqrt s}{(x+s)R_s}\dd s>0,\qquad
 \rho_{A,B}(x)=\frac{\sqrt{(b-x)(x-a)}}{2\pi}h_{A,B}(x).
\end{equation}
The density has mass one, positive endpoint coefficients, and its effective
potential derivative off the support is
$\sqrt{(x-a)(x-b)}h_{A,B}(x)$, with the square root negative below $a$.
This proves the strict exterior gap on every closed interval outside the
support. The functions extend holomorphically near every compact positive
interval. These are precisely the offcriticality and exterior properties
used below. Variation under $x\mapsto tx$ at its minimum gives
$B\int\sqrt x\dd\rho_{A,B}=2(A+1)$; hence
\begin{equation}\tag{AG7}
 \int\sqrt x\dd\rho=6/\pi.
\end{equation}
Indeed the energy of the pushed-forward measure is
$B\sqrt t\int\sqrt x\dd\rho-A\log t-A\int\log x\dd\rho
-\log t-\iint\log|x-y|\dd\rho\dd\rho$; its derivative at $t=1$ is zero.

Here is an explicit derivation of the required new moment, including its sign.
Set $r_0=\sqrt{ab}$, $m=(a+b)/2$, and
$R(z)=\sqrt{(z-a)(z-b)}$ with $R(z)/z\to1$. EIQ's elementary transform is
\[
 H_s(z)=1-\frac{R_s+R(z)}{z+s},\qquad
 G(z)=\frac{A}{2r_0}H_0(z)-\frac B{4\pi}
       \int_0^\infty\frac{H_s(z)}{\sqrt s R_s}\dd s.
\]
Since $R(0)=-r_0$, $R'(0)=m/r_0$, and
$\int_0^\infty (\sqrt sR_s)^{-1}\dd s=2\pi A/(Br_0)$,
evaluation at zero, followed by an integration by parts, gives
\begin{align}\tag{AG8}
 \int x^{-1}\dd\rho_{A,B}
 &=\frac{A(a+b)}{4ab}
   -\frac B{4\pi}\int_0^\infty
           \frac{1-r_0/R_s}{s^{3/2}}\dd s\\
 &=\frac{A(a+b)}{4ab}
   -\frac{Br_0}{4\pi}\int_0^\infty
         \frac{a+b+2s}{\sqrt s[(a+s)(b+s)]^{3/2}}\dd s.\notag
\end{align}
The boundary terms vanish: $1-r_0/R_s=O(s)$ at zero and is bounded at infinity.
Differentiating the absolutely convergent identity
$\int (\sqrt sR_s)^{-1}\dd s=2K(\kappa)/v$ with the operator
$-2(\partial_a+\partial_b)$ evaluates the last integral as $2E(\kappa)/(u^2v)$.
This derivative also follows directly from the defining elliptic integrals,
and every differentiated integral is integrable at both ends. Therefore
\begin{equation}\tag{AG9}
 M=\frac{u^2+v^2-4uv}{2u^2v^2}>0.
\end{equation}
Its positivity follows from the original positive probability measure.

\section{Full-form bounds on both complementary regions}
Write $r_k=k\sqrt{\delta^2+\gamma^2}/q$, so $r_k\to0$.
Pair opposite roots before estimating. The normalized factor
$E_q(x)=|Q(q\sqrt x)|^2/q^{2q}$ satisfies for $x>r_k^2$
\begin{equation}\tag{AG10}
 (x-r_k^2)^q\le E_q(x)\le(x+r_k^2)^q;
 \quad E_q(x)\le(x+r_k^2)^q\quad(x\ge0).
\end{equation}
Stirling's formula and positivity on compact real intervals give fixed
$c_\sigma,C_\sigma'>0$ such that, for all $y\ge0$,
\[
 c_\sigma(1+y)^{-1/2}e^{-\alpha y}
 \le\sigma(y)\le C_\sigma'(1+y)^{-1/2}e^{-\alpha y}.
\]
For $\varepsilon=0,1$ use the exact rescaled density
$g_{q,\varepsilon}(x)=q^{1/2}x^{\varepsilon-1/2}E_q(x)\sigma(q\sqrt x)$.
Every relevant polynomial has degree $d\le q/2$. Its shifted Legendre
expansion on $[1,2]$ proves the elementary pointwise bounds
\begin{equation}\tag{AG11}
 |P(x)|^2\le(d+1)^2
 \begin{cases}7^{2d},&0\le x\le1,\\(5x)^{2d},&x\ge2\end{cases}
 \int_1^2|P(t)|^2\dd t.
\end{equation}
For completeness, the normalized basis is
$\sqrt{2j+1}P_j(2x-3)$; Cauchy--Schwarz and
$|P_j(t)|\le(2|t|+1)^j$ outside $[-1,1]$ give (AG11), since
$\sum_{j=0}^d(2j+1)=(d+1)^2$. The polynomial bound follows from the
integral representation of $P_j$ or its real three-term recurrence.
On $[1,2]$, with $r_k^2\le1/2$ and $q\ge1$, the full density is at least
$c_0e^{-C_0q}$, where one may take
$C_0=\log2+\alpha\sqrt2$ and
$c_0=c_\sigma/[\sqrt2\sqrt{1+\sqrt2}]$.

Choose fixed $0<\epsilon<\min(1/2,u^2/2)$ with
$14\epsilon e^{C_0}<e^{-2}$. Impose the finite guard $r_k^2\le\epsilon$.
On $[0,\epsilon]$, (AG10)--(AG11) bound the integral of
$|P|^2g_{q,\varepsilon}$, relative to its $[1,2]$ integral, by
\begin{equation}\tag{AG12}
 \frac{C_\sigma'}{c_0}(d+1)^2
 (14\epsilon e^{C_0})^q
 \frac{\epsilon^{\varepsilon+1/4}}{\varepsilon+1/4}.
\end{equation}
The bound uses $q^{1/2}\sigma(q\sqrt x)\le
C_\sigma'x^{-1/4}e^{-\alpha q\sqrt x}$, valid for every $x>0$.
In particular it includes the entire region $q\sqrt x\le1$.

Choose $M_0>\max(2,2v^2,(8/\alpha)^2)$ such that
$\log10+C_0+2\log x\le\alpha\sqrt x/2$ for $x\ge M_0$.
It suffices to check this at $M_0$ and the displayed derivative guard.
For $x\ge M_0$, the same formulas bound the upper-tail ratio by
\begin{equation}\tag{AG13}
 \frac{C_\sigma'}{c_0}(d+1)^2
 \int_{M_0}^\infty x^{\varepsilon-3/4}
                  e^{-\alpha q\sqrt x/2}\dd x.
\end{equation}
Here $(5x)^{2d}(2x)^q\le(10x^2)^q$ and $d\le q/2$.
Substitution $z=\sqrt x$ shows that (AG13) is bounded by a fixed
polynomial in $q$ times $e^{-\alpha q\sqrt{M_0}/4}$.
Thus the sum $b_q$ of (AG12)--(AG13) is $O_h(q^2e^{-c_1q})$ for a
fixed $c_1>0$. These are inequalities on every vector before minimization:
\begin{equation}\tag{AG14}
 H_{[\epsilon,M_0]}\preceq H_{(0,\infty)}
 \preceq(1+b_q)H_{[\epsilon,M_0]},\qquad
 0\le\log\frac{\det H_{(0,\infty)}}{\det H_{[\epsilon,M_0]}}
 \le n\log(1+b_q).
\end{equation}
The identical argument with $E_q=x^q$ gives the comparison for the power
reference. No polynomial root, minimum or complementary integral is omitted.

\section{Exact fixed-potential block and a bounded-parameter lemma}
On the fixed interval, the real part of Stirling's expansion gives
\begin{equation}\tag{AG15}
 \sigma(q\sqrt x)=\sqrt2\,q^{-1/2}x^{-1/4}
 e^{-\pi q\sqrt x/2}(1+O(q^{-2})),
\end{equation}
uniformly in $x$. Its inverse-$y$ contribution has zero real part;
the next term is $O(y^{-2})$. This relative bound passes through every
Gram in (AG14), costing $O(nq^{-2})$ in its logarithmic determinant.
The finite root sums give
\begin{equation}\tag{AG16}
 \sum\omega_{ab}^2=\frac{qk(k+2)}3(\gamma^2-\delta^2),\qquad
 \log\frac{E_q(x)}{x^q}=-\frac{\tau_k}{x}+r_q(x),
 \quad \tau_k=\frac{\eta_hk(k+2)}q,\quad
 \eta_h=\frac{\gamma^2-\delta^2}3.
\end{equation}
All sums include multiplicity. Indeed
$\sum_{a=0}^k(2a-k)=0$ and
$\sum_{a=0}^k(2a-k)^2=(k+1)k(k+2)/3$.
In a fixed complex neighborhood of $[\epsilon,M_0]$ the logarithmic series
converges once $r_k^2/|x|\le1/2$; opposite roots remove odd terms and the
sum of the fourth and higher terms is bounded by
$C_h k^4/q^3$. Hence $\|r_q\|_\infty\le C_hk^4/q^3$.
Replacing $e^{r_q}$ in a complete determinant costs at most
$n\|r_q\|_\infty=O_h(k^4/q^2)=O_h(1)$, since $q\ge(k+1)^2$.

Let $q=2N$, $n=N+s$, and use $(s,\varepsilon)=(0,0),(1,0),(0,1)$.
After its exact powers of $q$ and $\sqrt2$ are extracted, the block density is
\begin{equation}\tag{AG17}
 x^{q+\varepsilon-3/4}e^{-\pi q\sqrt x/2-\tau_k/x}
 =e^{-nV(x)+W_{s,\varepsilon,\tau_k}(x)},\qquad
 W_{s,\varepsilon,t}=(\varepsilon-2s-3/4)\log x+\pi s\sqrt x-t/x.
\end{equation}
This equality, not a limiting replacement, fixes every parity increment.
The potential $V$ is now independent of $n$ and $k$.

Borot--Guionnet's soft-edge expansion and analytic linear-statistic formula
\cite[Propositions 1.2 and 5.2, equations (1-3)--(1-6)]{BG}, applied to
(AG6), give, for fixed analytic real $W$ on this interval,
\begin{equation}\tag{AG18}
 \log D_n[e^{-nV+W}]
 =n^2F+n\log(2\pi)+n\int W\dd\rho+O_W(\log n).
\end{equation}
Here $D_n$ is the monic Gram determinant. Its Vandermonde integral is the
partition function divided by $n!$. At ensemble index two the potential's
order-$n$ correction vanishes. The Gaussian reference is
$(2\pi)^{n/2}n^{-n^2/2}\prod_{j=0}^{n-1}j!$; Stirling summation gives
$-3n^2/4+n\log(2\pi)+O(\log n)$. This fixes the stated mass and factorial.
For the ratio with analytic factor the error after its order-$n$ term is
$O_W(1)$ by the cited linear-statistic formula. The support lies strictly
inside the actual integration interval, the potential is analytic there,
and (AG6) verifies its positive density and exterior gap; analyticity at
$x=0$ is not used.

We need no unproved uniformity in the varying $\tau_k$. Choose
$T>1+|\eta_h|$ and fix one of the three $W_{s,\varepsilon,0}$.
The function
\[
 f_n(t)=\log\frac{D_n[e^{-nV+W_{s,\varepsilon,t}}]}{D_n[e^{-nV}]}
                  -n\int W_{s,\varepsilon,t}\dd\rho
\]
is convex in $t$, by differentiating the positive squared-Vandermonde
integral twice (its second derivative is a variance).
Formula (AG18)'s ratio statement bounds $|f_n(-2T)|,|f_n(0)|,|f_n(2T)|$
by a common constant $C$. Convexity bounds $f_n(t)\le C$ on $[-T,T]$.
For $0\le t\le T$, write $0=\lambda t+(1-\lambda)(-2T)$,
$\lambda=2T/(t+2T)\ge2/3$; then $f_n(t)\ge-2C$.
The negative half is identical. This proves the uniform ratio error $O_h(1)$
for every actual $\tau_k$, since $|\tau_k|\le|\eta_h|$.
It includes fixed higher multiplicities without a fractional-power potential
expansion or a differentiated asymptotic remainder.

\section{Evaluate every block, source norm and root term}
The even block consists of $1,y^2,\ldots,y^{2n-2}$, and the odd block of
$y,y^3,\ldots,y^{2n-1}$. The substitution $y=q\sqrt x$, including both
halflines, contributes $q x^{-1/2}\dd x$.
Factoring the row and column powers and using (AG7), (AG14)--(AG18) proves
\begin{align}\tag{AG19}
 \log H_n^{q+\varepsilon}
 ={}&[2(q+\varepsilon)n+2n^2-\tfrac32n]\log q
       +\tfrac n2\log2+n^2F+n\log(2\pi)\\
 &+n[(\varepsilon-2s-\tfrac34)L+6s-\tau_kM]
       +O_h(\log q).\notag
\end{align}
The complete parity identity is
\begin{equation}\tag{AG20}
 \log R_q+\log R_{q+1}
 =\log H_N^q+\log H_{N+1}^q+2\log H_N^{q+1}.
\end{equation}
Their dimensions sum to $2q+1$. Their squared dimensions sum to
$q^2+q+1$; their $s$-weighted dimensions sum to $q/2+1$;
and their $\varepsilon$-weighted dimensions sum to $q$.
Keeping these exact sums until the end gives
\begin{align}\tag{AG21}
 \log R_q+\log R_{q+1}
 ={}&(6q^2+3q+\tfrac12)\log q+Fq^2\\
 &+q[F+3-\tfrac32L+3\log2+2\log\pi]
       -2\eta_h M k(k+2)+O_h(\log q).\notag
\end{align}
The extra $-\tau_kM$ from $2q+1$ is bounded and included in the displayed
error. Every constant dropped here is bounded independently of $k$.

For $R_1$ the reference integral is exactly
$2\sqrt2\,\Gamma(2q+1/2)/\alpha^{2q+1/2}$.
On $[\epsilon,M_0]$ the full root multiplier has logarithm bounded by
$|\tau_k|/\epsilon+\|r_q\|_\infty$, and (AG15) has relative error $O(q^{-2})$.
Both references satisfy (AG14), so their full integral ratio has bounded
logarithm. Stirling now gives
\begin{equation}\tag{AG22}
 \log R_1=2q\log q+(4\log2-2\log\pi-2)q+O_h(\log q).
\end{equation}
For the source norms, (AG3) yields uniformly for $q\le n\le2q$
$\log\gamma_n=2n\log n-2n+\tfrac12\log n+\tfrac32\log2+\log\pi+O(n^{-1})$.
The coefficients in (AG4) are twice the trapezoidal weights. Integrating
$2x\log x-2x$ and $\tfrac12\log x$ from $q$ to $2q$, and bounding the
trapezoidal remainder by its second derivative integral, gives
\begin{equation}\tag{AG23}
 \sum_{n=q}^{2q}c_{n-q}\log\gamma_n
 =6q^2\log q+(8\log2-9)q^2
  +q\log q+(5\log2+2\log\pi-1)q+O(\log q).
\end{equation}
For example $2\int_q^{2q}(2x\log x-2x)\dd x
=6q^2\log q+(8\log2-9)q^2$ and
$\int_q^{2q}\log x\dd x=q\log q+(2\log2-1)q$;
the constant part contributes $(3\log2+2\log\pi)q$.

\begin{theorem}[Absolute Gamma value on every fixed multiplicity]
With all original definitions retained, put
\begin{equation}\tag{AG24}
 C_B=9-8\log2+F,\qquad
 C_\sigma=F+6-\tfrac32L-6\log2+2\log\pi.
\end{equation}
Then
\begin{equation}\tag{AG25}
 \boxed{\BB_k[\sigma]=C_Bq^2+C_\sigma q
       -2\eta_hM k(k+2)+O_h(\log q).}
\end{equation}
\end{theorem}
\begin{proof}
Substitute (AG21)--(AG23) into the exact identity (AG4).
The $q\log q$ coefficient is $3-2-1=0$.
The remaining order-$q$ coefficient is
$F+3-\tfrac32L+3\log2+2\log\pi
-(4\log2-2\log\pi-2)-(5\log2+2\log\pi-1)$,
which is exactly $C_\sigma$. All preceding bounds are uniform in the
bounded $\tau_k$ and use $k^4/q^2\le1$; they apply to every fixed $m_0$.
\end{proof}

\section{Backwards and forwards through the original action}
Retain $w_h=|(2\xi/h)(1/2+iy)|^2/(2\pi)$, $m_k=w_h^{*k}$,
$I_h=\int_0^{\pi/2}[(\log M_h)'(\theta)]^2\dd\theta$,
$d=4m_0-2$, $c_\zeta=2\gamma_{\rm E}-\log(2\pi)$ and
$C_\Gamma=2L-8\log2+4$. The already proved SR1--48 receiver gives
\begin{equation}\tag{AG26}
 \BB_k[m_k]-\BB_k[\sigma]
 =-dC_\Gamma q+\frac{\log2}{\pi}I_hk^2
  +\frac{C_\Gamma q}{2(\log q+c_\zeta)}+R_{h,k},
 \qquad R_{h,k}=o_h(q/\log q).
\end{equation}
Its full finite remainder, annulus guards, inner and outer integrals and
the eight positive Grams are retained in
\texttt{SLOW\_TAIL\_AND\_NONLINEAR\_RECEIVER.tex}; its result is used
without differentiating $R_{h,k}$ or assigning its constituent signed pieces.
Adding (AG25) to that exact scalar difference proves
\begin{align}\tag{AG27}
 \BB_k[m_k]={}&C_Bq^2+(C_\sigma-dC_\Gamma)q
 +\frac{\log2}{\pi}I_hk^2-2\eta_hMk(k+2)\\
 &+\frac{C_\Gamma q}{2(\log q+c_\zeta)}
 +R_{h,k}+O_h(\log q).\notag
\end{align}
This evaluates the absolute baseline through the $q/\log q$ scale.
It strengthens the incoming $o(q)$ statement; no new arithmetic estimate
is presumed. The simple-packet root term contributes at order $q$.
At higher fixed multiplicity its displayed $k^2$ value is retained even
though the arithmetic error does not yet resolve all terms of that size.

The original CAI15 identity is
$J_k^{\rm action}=\BB_k[m_k]+\mathcal W_k^{\rm ar}-\mathcal P_--\mathcal P_+$.
Its finite guards and penalties give (CAI29--35)
\begin{gather}\tag{AG28}
 -E_{\rm hi}\le\mathcal W_k^{\rm ar}-\mathcal W_q^\sigma\le E_{\rm lo},
 \quad0\le\mathcal P_-+\mathcal P_+\le\Pi_{h,k},\\
 \mathcal W_q^\sigma=2\log\frac{(4q)!}{(2q)!}-4q\log2
       +\log\frac{2q-1}{2(4q-1)},\qquad
 E_{\rm hi}+E_{\rm lo}+\Pi_{h,k}=O_h(k+\log q).\notag
\end{gather}
In particular its exact factorial window is
$4q\log q+(8\log2-4)q-\log2+O(q^{-1})$.
Returning (AG27) through (AG28) proves
\begin{equation}\tag{AG29}
 \boxed{\begin{aligned}J_k^{\rm action}={}&C_Bq^2+4q\log q
  +[C_\sigma+8\log2-4-dC_\Gamma]q\\
 &+\frac{\log2}{\pi}I_hk^2-2\eta_hMk(k+2)
  +\frac{C_\Gamma q}{2(\log q+c_\zeta)}+o_h(q/\log q).
 \end{aligned}}
\end{equation}
Indeed $(k+\log q)/(q/\log q)\to0$ on the original sequence. At finite $k$
the exact interval (AG28) and the complete SR error remain the bounds,
rather than an assigned coefficient for any separate penalty or native loss.
The constant identity
\begin{equation}\tag{AG30}
 C_\sigma=C_B-\tfrac34 C_\Gamma+2\log(\pi/4)
\end{equation}
follows by substituting the original definitions; no mass is changed.
For the original benchmark $L_{h,k}=\delta q(k+1)/2$, subtracting its
actual logarithm gives the further receiver
\begin{align}\tag{AG31}
 J_k^{\rm action}-4q\log L_{h,k}
 ={}&C_Bq^2-4q\log(k+1)
 +[C_\sigma+8\log2-4-dC_\Gamma-4\log(\delta/2)]q\\
 &+\frac{\log2}{\pi}I_hk^2-2\eta_hMk(k+2)
 +\frac{C_\Gamma q}{2(\log q+c_\zeta)}+o_h(q/\log q).\notag
\end{align}
The full observation kernel still enters through its original inclusion
$I$ and quotient $\Lambda$: $H_{K,N}=I^*G_NI$ and
$Q_{B,N}=(\Lambda G_N^{-1}\Lambda^*)^{-1}$.
For any fixed section $S_*$ of $\Lambda$ the exact map
$[I,G_N^{-1}\Lambda^*Q_{B,N}]$ has determinant equal to $[I,S_*]$ and
orthogonal columns between its two blocks. Thus
\begin{equation}\tag{AG32}
 \det H_{K,N}\det Q_{B,N}=|\det[I,S_*]|^2\det G_N,
 \qquad F_K+F_B=\BB_k[m_k].
\end{equation}
Equation (AG27) replaces the previously unevaluated total in this equality;
it does not allocate that total to either original block. The independent
proper-source return and the original projected-current response ratios
retain their own exact maps and remain subsequent calculations.

\begin{thebibliography}{9}
\bibitem{BG} G. Borot and A. Guionnet,
\emph{Asymptotic expansion of beta matrix models in the one-cut regime},
Communications in Mathematical Physics 317 (2013), 447--483;
\href{https://arxiv.org/abs/1107.1167}{arXiv:1107.1167v3},
Hypothesis 1.1, Propositions 1.2 and 5.2, Sections 5.1--5.3.
\bibitem{DLMF} F. W. J. Olver et al., editors,
\emph{NIST Digital Library of Mathematical Functions},
\href{https://dlmf.nist.gov/18.23.E7}{equation 18.23.7},
Meixner--Pollaczek generating function; and
\href{https://dlmf.nist.gov/5.11}{Section 5.11}, Gamma asymptotics.
\end{thebibliography}
\end{document}
